\begin{lemma}
\label{lemma-good-right-resolution-gives-qis}
Let $M^\bullet$ be a complex of abelian groups. Let
$$
0 \to M^\bullet \to A_0^\bullet \to A_1^\bullet \to A_2^\bullet \to \ldots
$$
be an exact complex of complexes of abelian groups
such that for all $p \in \mathbf{Z}$ the complexes
$$
0 \to
\Coker(d_{M^\bullet}^p) \to
\Coker(d_{A_0^\bullet}^p) \to
\Coker(d_{A_1^\bullet}^p) \to
\Coker(d_{A_2^\bullet}^p) \to \ldots
$$
are exact as well. Set $A^{p, q} = A_p^q$ to obtain a double
complex. Let $\text{Tot}_\pi(A^{\bullet, \bullet})$ be the
product total complex associated to the double complex
(see proof). Then the map
$M^\bullet \to \text{Tot}_\pi(A^{\bullet, \bullet})$
induced by $M^\bullet \to A_0^\bullet$ is a quasi-isomorphism.
\end{lemma}

\begin{proof}
Abbreviating $T^\bullet = \text{Tot}_\pi(A^{\bullet, \bullet})$
we define
$$
T^n = \prod\nolimits_{p + q = n} A^{p, q} =
\prod\nolimits_{p + q = n} A_p^q
\quad\text{with}\quad
\text{d}_{T^\bullet}^n =
\prod\nolimits_{n = p + q} (f_p^q + (-1)^pd_{A_p^\bullet}^q)
$$
where $f_p^\bullet : A_p^\bullet \to A_{p + 1}^\bullet$
are the maps of complexes in the lemma.

\medskip\noindent
We will show that $H^0(M^\bullet) \to H^0(T^\bullet)$ is an isomorphism.
The same argument works for other degrees.
Let $x \in \Ker(\text{d}_{T^\bullet}^0)$ represent $\xi \in H^0(T^\bullet)$.
Write $x = (x_i)$ with $x_i \in A_i^{-i}$.
Note that $x_0$ maps to zero in $\Coker(A_1^{-1} \to A_1^0)$.
Hence we see that $x_0 = m_0 + d_{A_0^\bullet}^{-1}(y)$ for some
$m_0 \in M^0$ and $y \in A_0^{-1}$.
Then $d_{M^\bullet}(m_0) = 0$ because $\text{d}_{A_0^\bullet}(x_0) = 0$
as $\text{d}_{T^\bullet}(x) = 0$.
Thus, replacing $\xi$ by something in the image of
$H^0(M^\bullet) \to H^0(T^\bullet)$ we may assume that $x_0$
is in $\Im(d^{-1}_{A_0^\bullet})$.

\medskip\noindent
Assume $x_0 \in \Im(d^{-1}_{A_0^\bullet})$. We claim that in this
case $\xi = 0$. To prove this we find, by induction on $n$ elements
$y_0, y_1, \ldots, y_n$ with $y_i \in A_i^{-i - 1}$ such that
$x_0 = \text{d}_{A_0}^{-1}(y_0)$ and
$x_j = f_{j - 1}^{-j}(y_{j - 1}) + (-1)^j d^{-j - 1}_{A_j^\bullet}(y_j)$
for $j = 1, \ldots, n$. This is clear for $n = 0$. Proof of induction step:
suppose we have found $y_0, \ldots, y_{n - 1}$. Then
$w_n = x_n - f_{n - 1}^{-n}(y_{n - 1})$ is in the kernel of
$d^{-n}_{A_n^\bullet}$ and maps to zero in $H^{-n}(A_{n + 1}^\bullet)$
(because it maps to an element which is the boundary
of $x_{n + 1}$ up to sign). Exactly as in the proof of
Lemma \ref{lemma-good-resolution-gives-qis}
the assumptions of the lemma imply that
$$
0 \to
H^p(M^\bullet) \to
H^p(A_0^\bullet) \to
H^p(A_1^\bullet) \to
H^p(A_2^\bullet) \to \ldots
$$
is exact for all $p \in \mathbf{Z}$. Thus after changing $y_{n - 1}$
by an element in $\Ker(d^{n - 1}_{A_{n - 1}^\bullet})$ we may assume
that $w_n$ maps to zero in $H^{-n}(A_n^\bullet)$. This means we
can find $y_n$ as desired. Observe that this procedure does not
change $y_0, \ldots, y_{n - 2}$. Hence continuing ad infinitum
we find an element $y = (y_i)$ in $T^{n - 1}$ with $d_{T^\bullet}(y) = \xi$.
This shows that $H^0(M^\bullet) \to H^0(T^\bullet)$ is surjective.

\medskip\noindent
Suppose that $m_0 \in \Ker(d^0_{M^\bullet})$ maps to zero in $H^0(T^\bullet)$.
Say it maps to the differential applied to $y = (y_i) \in T^{-1}$ .
Then $y_0 \in A_0^{-1}$ maps to zero in $\Coker(d^{-2}_{A_1^\bullet})$.
By assumption this means that $y_0 \bmod \Im(d^{-2}_{A_0^\bullet})$
is the image of some $z \in M^{-1}$. It follows that
$m_0 = d^{-1}_{M^\bullet}(z)$. This proves injectivity and the proof is
complete.
\end{proof}